% =============================================================
%  Relativistic modification of Chapter VII, §71
%  Narrow gap Couette solutions — relativistic formulation
%
%  Original: Chandrasekhar, §71 "The solution for the case of a
%            narrow gap when the marginal state is stationary"
%  This file extends the analysis to a relativistic fluid using
%  the BDNK (Bemfica–Disconzi–Noronha–Kovtun) causal framework.
%
%  Key modification: the inertia ρ is replaced by (ε+p)/c² in
%  the Taylor number definition and the perturbation equations,
%  yielding Ta_rel and modified critical Taylor numbers.
% =============================================================

\section{Narrow gap Couette solutions for the marginal stationary
state --- relativistic}\label{sec:rel-71}

In \S\,\ref{sec:7-71} we derived the equations governing the onset of
instability in the narrow-gap limit and obtained the critical Taylor
number~$T_{c}$ as a function of~$\mu = \Omega_{2}/\Omega_{1}$.  We now
extend these results to a relativistic fluid in which the effective
inertia is $\enthalpy/c^{2} = (\edensity + p)/c^{2}$ rather than the
rest-mass density~$\rdensity$ alone.  As in the preceding relativistic
sections we keep~$c$ explicit so that every correction can be traced to
its physical origin; the non-relativistic limit is recovered by
letting $c \to \infty$.

Throughout this section the notation of \S\,\ref{sec:7-71} is retained:
$d = R_{2} - R_{1}$ is the gap width,
$\zeta = (r - R_{1})/d$,
$a = k\,d$ is the dimensionless wave number, and
$D = \dd/\dd\zeta$.

%--------------------------------------------------------------
\subsection{Relativistic perturbation equations in the narrow gap}
\label{sec:rel-71-equations}
%--------------------------------------------------------------

For a relativistic perfect fluid in circular Couette flow the
energy-momentum tensor is
\begin{equation}
\label{eq:rel-7-191}
\emt = (\edensity + p)\,\frac{u^{\mu}\,u^{\nu}}{c^{2}} + p\,\metric.
\end{equation}
In the narrow-gap approximation the curvature terms are negligible
and the background angular velocity profile retains the classical form
\begin{equation}
\label{eq:rel-7-192}
\Omega(\zeta)
  = \Omega_{1}\bigl[1 - (1-\mu)\,\zeta\bigr],
\end{equation}
but the linearised momentum equations acquire the replacement
$\rdensity \to \enthalpy/c^{2}$ in every inertial term.
Defining the \emph{relativistic kinematic viscosity}
\begin{equation}
\label{eq:rel-7-193}
\nu_{\mathrm{rel}}
  = \frac{\shearvisc}{\enthalpy/c^{2}}
  = \frac{\shearvisc\,c^{2}}{\edensity + p},
\end{equation}
where $\shearvisc$ is the dynamic shear viscosity, the narrow-gap
perturbation equations~\eqref{eq:7-196}--\eqref{eq:7-197} become
\begin{equation}
\label{eq:rel-7-196}
\bigl(D^{2} - a^{2} - \sigma_{\mathrm{rel}}\bigr)
\bigl(D^{2} - a^{2}\bigr)\,u
  = (1 + \alpha\,\zeta)\,v,
\end{equation}
\begin{equation}
\label{eq:rel-7-197}
\bigl(D^{2} - a^{2} - \sigma_{\mathrm{rel}}\bigr)\,v
  = -T_{\mathrm{rel}}\,a^{2}\,u,
\end{equation}
with
\begin{equation}
\label{eq:rel-7-198}
\sigma_{\mathrm{rel}} = \frac{p\,d^{2}}{\nu_{\mathrm{rel}}}
\end{equation}
and the \emph{relativistic Taylor number}
\begin{equation}
\label{eq:rel-7-199}
\boxed{\;
T_{\mathrm{rel}}
  = \frac{4\,A\,\Omega_{1}}{\nu_{\mathrm{rel}}^{\,2}}\,d^{4}
  = \frac{4\,A\,\Omega_{1}\,(\edensity+p)^{2}}
         {\shearvisc^{2}\,c^{4}}\,d^{4}.
\;}
\end{equation}
Comparing with the classical Taylor number~\eqref{eq:7-198},
\begin{equation}
\label{eq:rel-7-200}
T_{\mathrm{rel}}
  = T_{\mathrm{cl}}
    \left(\frac{\edensity + p}{\rdensity\,c^{2}}\right)^{\!2}
  = T_{\mathrm{cl}}
    \left(1 + \frac{\edensity + p - \rdensity c^{2}}
                    {\rdensity\,c^{2}}\right)^{\!2}.
\end{equation}
For a non-relativistic fluid $\edensity + p \approx \rdensity c^{2}$
and $T_{\mathrm{rel}} \to T_{\mathrm{cl}}$ as expected.  The
correction is always \emph{positive}: thermal and pressure
contributions to the inertia raise the effective Taylor number,
making the flow \emph{more unstable} at a given angular velocity.

\begin{relcorrection}
The classical Taylor number
$T = 4A\Omega_{1}d^{4}/\nu^{2}$
is replaced by
$T_{\mathrm{rel}} = 4A\Omega_{1}d^{4}/\nu_{\mathrm{rel}}^{2}$,
where $\nu_{\mathrm{rel}} = \shearvisc\,c^{2}/(\edensity+p)$.
The ratio $T_{\mathrm{rel}}/T_{\mathrm{cl}}
= [(\edensity+p)/(\rdensity c^{2})]^{2}$ quantifies the
relativistic enhancement of centrifugal instability.
\end{relcorrection}

The parameter $\alpha = -(1-\mu)$ is \emph{unchanged} because it
characterises the angular-velocity profile, not the fluid inertia.

The boundary conditions remain
\begin{equation}
\label{eq:rel-7-201}
u = Du = v = 0
\quad\text{for } \zeta = 0 \text{ and } 1.
\end{equation}

%--------------------------------------------------------------
\subsection{Characteristic value problem: the relativistic
critical Taylor number}\label{sec:rel-71-eigenvalue}
%--------------------------------------------------------------

Setting $\sigma_{\mathrm{rel}} = 0$ (stationary marginal state)
the equations to be solved are identical in \emph{form} to the
classical equations~\eqref{eq:7-201}--\eqref{eq:7-202}, with
$T_{\mathrm{cl}}$ replaced everywhere by~$T_{\mathrm{rel}}$:
\begin{equation}
\label{eq:rel-7-202}
(D^{2}-a^{2})^{2}\,u = (1+\alpha\,\zeta)\,v,
\end{equation}
\begin{equation}
\label{eq:rel-7-203}
(D^{2}-a^{2})\,v = -T_{\mathrm{rel}}\,a^{2}\,u.
\end{equation}
The sine-series expansion $v = \sum_{m} C_{m}\sin m\pi\zeta$ and
the subsequent Galerkin procedure of \S\,\ref{sec:7-71}\,(a) carry
over verbatim; the secular determinant~\eqref{eq:7-213} now
determines $T_{\mathrm{rel}}$ instead of~$T_{\mathrm{cl}}$.

\paragraph{First approximation.}
The $(1,1)$-element of the secular matrix yields the relativistic
analogue of equation~\eqref{eq:7-215}:
\begin{equation}
\label{eq:rel-7-215}
T_{\mathrm{rel}}
  = \frac{2}{2+\alpha}\,
    \frac{(\pi^{2}+a^{2})^{3}}{a^{2}
      \bigl[1
       - 16a^{2}\pi^{2}\cosh^{2}\!\tfrac{1}{2}a
         /(\pi^{2}+a^{2})^{2}(\sinh a + a)\bigr]}.
\end{equation}
Since the right-hand side is identical to the classical expression
(differing only in the interpretation of~$T$), the \emph{critical
relativistic Taylor number} is
\begin{equation}
\label{eq:rel-7-216}
\boxed{\;
T_{c}^{\mathrm{rel}}
  = \frac{3430}{1+\mu}
    \left(\frac{\edensity+p}{\rdensity\,c^{2}}\right)^{\!2},
  \qquad
  a_{\min} = 3{\cdot}12.
\;}
\end{equation}
Equivalently, in terms of the classical critical Taylor number
$T_{c}^{\mathrm{cl}} = 3430/(1+\mu)$, we have
\begin{equation}
\label{eq:rel-7-217}
\frac{T_{c}^{\mathrm{rel}}}{T_{c}^{\mathrm{cl}}}
  = \left(\frac{\edensity+p}{\rdensity\,c^{2}}\right)^{\!2}
  = 1 + \frac{2(\edensity + p - \rdensity c^{2})}{\rdensity c^{2}}
      + \order{c^{-4}}.
\end{equation}

\begin{relcorrection}
The critical Taylor number for the onset of instability is
\emph{raised} by the factor $[(\edensity+p)/(\rdensity c^{2})]^{2}$
relative to the classical value.  The wave number at onset is
unaffected to leading relativistic order.
\end{relcorrection}

%--------------------------------------------------------------
\subsection{Numerical results: $T_{c}^{\mathrm{rel}}$ versus
$T_{c}^{\mathrm{cl}}$}\label{sec:rel-71-numerical}
%--------------------------------------------------------------

Because the secular determinant~\eqref{eq:7-213} is structurally
unchanged—only the meaning of~$T$ differs—all classical numerical
results of Table~\ref{tab:XXXII} and Table~\ref{tab:XXXIII} carry
over to the relativistic case upon multiplication by the
enthalpy-inertia factor.  We define the \emph{relativistic
correction parameter}
\begin{equation}
\label{eq:rel-7-218}
\xi \;=\; \frac{\edensity + p}{\rdensity\,c^{2}} - 1
      \;=\; \frac{\edensity + p - \rdensity c^{2}}{\rdensity\,c^{2}},
\end{equation}
so that
\begin{equation}
\label{eq:rel-7-219}
T_{c}^{\mathrm{rel}}
  = T_{c}^{\mathrm{cl}}\,(1+\xi)^{2}
  \approx T_{c}^{\mathrm{cl}}
          \bigl(1 + 2\xi + \xi^{2}\bigr).
\end{equation}

\begin{table}[htbp]
\centering
\caption{Ratio of relativistic to classical critical Taylor numbers
for representative values of the correction parameter~$\xi$.}
\label{tab:rel-XXXII}
\begin{tabular}{ccl}
\hline
$\xi$ & $(1+\xi)^{2}$ &
  Typical physical regime \\
\hline
$10^{-10}$ & $\approx 1$ &
  Terrestrial laboratory fluid \\
$10^{-5}$ & $1 + 2\times 10^{-5}$ &
  White-dwarf interior \\
$10^{-2}$ & $1{\cdot}0201$ &
  Neutron-star outer core \\
$0{\cdot}1$ & $1{\cdot}21$ &
  Hot neutron-star interior \\
$0{\cdot}5$ & $2{\cdot}25$ &
  Quark--gluon plasma \\
$1{\cdot}0$ & $4{\cdot}0$ &
  Ultra-relativistic limit ($p \sim \edensity$) \\
\hline
\end{tabular}
\end{table}

The table shows that the relativistic correction is negligible for
terrestrial fluids but can be substantial in astrophysical settings
where~$p/\edensity$ is not small.

The variation with~$\mu$ remains identical to the classical case at
every order of approximation.  In particular the successive Galerkin
approximations converge at the same rate because the matrix
structure of the secular equation is unchanged; only the diagonal
and off-diagonal entries are rescaled uniformly by the
enthalpy-inertia factor.  The coefficients $\mathscr{C}_{m}$ listed
in Table~\ref{tab:XXXIII} are therefore also unchanged.

%--------------------------------------------------------------
\subsection{Alternative method of solution}\label{sec:rel-71-alternative}
%--------------------------------------------------------------

The alternative formulation of \S\,\ref{sec:7-71}\,(c), in which
$u$ is eliminated between equations~\eqref{eq:rel-7-202}
and~\eqref{eq:rel-7-203} to obtain a sixth-order equation for~$v$
alone, carries over to the relativistic case with no modification
beyond the substitution $T_{\mathrm{cl}} \to T_{\mathrm{rel}}$.
The resulting equation is
\begin{equation}
\label{eq:rel-7-220}
(D^{2}-a^{2})^{3}\,v = -T_{\mathrm{rel}}\,a^{2}(1+\alpha\,\zeta)\,v,
\end{equation}
with boundary conditions
\begin{equation}
\label{eq:rel-7-221}
v = (D^{2}-a^{2})\,v = D(D^{2}-a^{2})\,v = 0
\quad\text{for } \zeta = 0 \text{ and } 1.
\end{equation}

The secular equation~\eqref{eq:7-228}, with~$T$ interpreted as
$T_{\mathrm{rel}}$, again yields numerical values that are
`practically indistinguishable' from those of the first method,
exactly as in the classical theory.  Explicitly, for
$\mu = -0{\cdot}5$ and $a = 3{\cdot}20$:
\begin{equation}
\label{eq:rel-7-230}
T_{\mathrm{rel}} = (1+\xi)^{2} \times
\begin{cases}
6417{\cdot}8 & \text{(6th-order method, 3rd approximation)},\\
6417{\cdot}1 & \text{(4th-order method, 3rd approximation)}.
\end{cases}
\end{equation}
The near-equality of the two methods persists because the
relativistic modification is a uniform multiplicative rescaling
that does not alter the matrix structure or the rate of convergence
of the Galerkin expansion.

%--------------------------------------------------------------
\subsection{Approximate solution for $\mu \to 1$}\label{sec:rel-71-approx}
%--------------------------------------------------------------

The perturbation expansion developed in
\S\,\ref{sec:7-71}\,(d)---which exploits the proximity of the
narrow-gap Couette problem to the B\'enard problem when
$\mu \to 1$---generalises in a straightforward manner.

Introducing the shifted coordinate $x = \zeta - \tfrac{1}{2}$ and
\begin{equation}
\label{eq:rel-7-234}
\lambda_{\mathrm{rel}}
  = \tfrac{1}{2}\,T_{\mathrm{rel}}\,a^{2}(1+\mu),
\qquad
\epsilon = -\frac{2(1-\mu)}{1+\mu},
\end{equation}
equation~\eqref{eq:rel-7-220} becomes
\begin{equation}
\label{eq:rel-7-235}
(D^{2}-a^{2})^{3}\,v
  = -\lambda_{\mathrm{rel}}\,(1+\epsilon\,x)\,v.
\end{equation}
This is \emph{identical in form} to the classical
equation~\eqref{eq:7-233}: only the characteristic value
$\lambda_{\mathrm{rel}}$ carries the relativistic content through
its dependence on~$T_{\mathrm{rel}}$.

The perturbation theory then proceeds exactly as before.  The
zero-order characteristic value is
$\lambda_{0} = R_{c}\,a^{2} = 1708\,a^{2}$ (for $a = 3{\cdot}117$),
the first-order correction vanishes ($\lambda^{(1)} = 0$) by the
parity argument, and the second-order correction is given by
equation~\eqref{eq:7-258}.  Consequently,
\begin{equation}
\label{eq:rel-7-263}
\boxed{\;
T_{c}^{\mathrm{rel}}
  = \frac{3416}{1+\mu}
    \left(\frac{\edensity+p}{\rdensity\,c^{2}}\right)^{\!2}
    \!\left\{1 - 7{\cdot}61 \times 10^{-2}
      \left(\frac{1-\mu}{1+\mu}\right)^{\!2}\right\}.
\;}
\end{equation}
The relativistic correction enters as a \emph{multiplicative
prefactor} that does not disturb the~$\mu$-dependent structure of
the perturbation series.  The reason is clear: the enthalpy-inertia
factor $(\edensity+p)/c^{2}$ is spatially uniform in the narrow-gap
limit and therefore commutes with the differential operators in the
eigenvalue problem.

\begin{relcorrection}
For $\mu \to 1$ the approximate formula
$T_{c}^{\mathrm{cl}} = 3416/(1+\mu)$
generalises to
$T_{c}^{\mathrm{rel}} = 3416\,(\edensity+p)^{2}/
[\rdensity^{2}c^{4}(1+\mu)]$,
with $\order{\epsilon^{2}}$ corrections identical to those in the
classical theory.  The first-order correction in~$\epsilon$
vanishes by the same symmetry argument as in the classical case.
\end{relcorrection}

%--------------------------------------------------------------
\subsection{Asymptotic behaviour for $(1-\mu) \to \infty$}
\label{sec:rel-71-asymptotic}
%--------------------------------------------------------------

In the regime of strong counter-rotation,
$(1-\mu) \to \infty$, the classical analysis of
\S\,\ref{sec:7-71}\,(e) established the asymptotic relations
\begin{equation}
\label{eq:rel-7-264}
T_{c}^{\mathrm{cl}} \to \tau\,(1-\mu)^{4},
\qquad
a(T_{c}) \to q\,(1-\mu),
\end{equation}
with $\tau \approx 1182$ and $q \approx 2{\cdot}035$.

\paragraph{Physical origin of the scaling.}
The instability is localised to a layer of thickness
$d_{0} = d/(1-\mu)$ adjoining the inner cylinder, where the
Rayleigh discriminant is violated.  Since the Taylor number scales
as the fourth power of the effective gap, the $(1-\mu)^{4}$
behaviour follows.

In the relativistic theory the same physical picture holds.  The
angular-velocity profile~\eqref{eq:rel-7-192} is unmodified and the
nodal surface remains at $\zeta = 1/(1-\mu)$.  The only change is
the replacement of~$\nu$ by~$\nu_{\mathrm{rel}}$ in the definition
of~$T$.  Hence
\begin{equation}
\label{eq:rel-7-267}
\boxed{\;
T_{c}^{\mathrm{rel}}
  \;\to\; \tau\,(1-\mu)^{4}
    \left(\frac{\edensity+p}{\rdensity\,c^{2}}\right)^{\!2}
  \quad\text{as } (1-\mu) \to \infty,
\;}
\end{equation}
with $\tau \approx 1182$ unchanged, and
\begin{equation}
\label{eq:rel-7-268}
a\bigl(T_{c}^{\mathrm{rel}}\bigr)
  \;\to\; q\,(1-\mu),
  \quad q \approx 2{\cdot}035.
\end{equation}
The wave number at onset is unaffected because it is determined by
the geometry of the unstable layer, not by the magnitude of the
inertia.

\paragraph{Analytical structure.}
With the substitution $x = a\,\zeta$,
equation~\eqref{eq:rel-7-220} becomes
\begin{equation}
\label{eq:rel-7-270}
(D^{2}-1)^{3}\,v
  = -\frac{T_{\mathrm{rel}}}{a^{4}}
    \!\left(1 - \frac{1-\mu}{a}\,x\right) v.
\end{equation}
In the limit $(1-\mu) \to \infty$, $a \to \infty$ with
$\gamma = (1-\mu)/a$ and
$\lambda_{\mathrm{rel}} = T_{\mathrm{rel}}/a^{4}$ held fixed, one
obtains the semi-infinite eigenvalue problem
\begin{equation}
\label{eq:rel-7-273}
(D^{2}-1)^{3}\,v = -\lambda_{\mathrm{rel}}\,(1-\gamma\,x)\,v,
\end{equation}
\begin{equation}
\label{eq:rel-7-274}
v = (D^{2}-1)\,v = D(D^{2}-1)\,v = 0
\quad\text{for } x = 0 \text{ and } \infty.
\end{equation}
The structure is identical to the classical
equations~\eqref{eq:7-273}--\eqref{eq:7-274}; hence $\tau$ and $q$
are the same universal constants as in the Newtonian theory.  The
enthalpy-inertia correction enters only through
$\lambda_{\mathrm{rel}} = T_{\mathrm{rel}}/a^{4}$ and therefore
factors out of the minimisation
$\tau = \min_{\gamma}\{\lambda_{\mathrm{rel}}/\gamma^{4}\}$.

\begin{relcorrection}
The asymptotic relations
$T_{c} \sim \tau(1-\mu)^{4}$ and
$a(T_{c}) \sim q(1-\mu)$ for $(1-\mu) \to \infty$ persist in
the relativistic theory with the same universal constants
$\tau \approx 1182$ and $q \approx 2{\cdot}035$.
The relativistic correction is an overall multiplicative factor
$[(\edensity+p)/(\rdensity c^{2})]^{2}$ on~$T_{c}$, reflecting
the enhanced inertia.  The wave number is unaffected because the
geometry of the unstable layer is independent of the equation of
state.
\end{relcorrection}

%--------------------------------------------------------------
\subsection{Summary and discussion}\label{sec:rel-71-summary}
%--------------------------------------------------------------

The relativistic generalisation of the narrow-gap Couette stability
problem is governed by a single modification: the replacement of the
rest-mass density~$\rdensity$ by the relativistic inertia
$(\edensity + p)/c^{2}$ in the definition of the Taylor number.
This replacement:

\begin{enumerate}
\item Multiplies all critical Taylor numbers by the universal
  factor $[(\edensity+p)/(\rdensity c^{2})]^{2}$.

\item Leaves the wave number at onset, the convergence of the
  Galerkin expansion, and the eigenfunction coefficients
  $\mathscr{C}_{m}$ unchanged.

\item Preserves the validity of both the first method (4th-order
  system for~$u,v$) and the alternative method (6th-order system
  for~$v$ alone) at equal rates of convergence.

\item Preserves the approximate formula
  $T_{c} = 3416\,(1+\xi)^{2}/(1+\mu)$ for $\mu \to 1$, including
  the vanishing of the first-order correction in~$\epsilon$.

\item Preserves the asymptotic scaling
  $T_{c} \sim \tau(1-\mu)^{4}(1+\xi)^{2}$ for
  $(1-\mu) \to \infty$, with the same universal constants as in
  the classical theory.
\end{enumerate}

The physical reason for this structural simplicity is that, in the
narrow-gap limit, the enthalpy-inertia factor is spatially uniform
and therefore acts as a constant rescaling of the effective
viscosity.  The eigenvalue problem is consequently \emph{isospectral}
to the classical one up to this overall factor.

For astrophysical applications---accretion discs, neutron-star
interiors, and the quark--gluon plasma---the correction can be
significant.  In the ultra-relativistic limit
$p \sim \edensity/3$ (radiation-dominated fluid) one has
$\xi \approx 1/3$, giving
$T_{c}^{\mathrm{rel}}/T_{c}^{\mathrm{cl}} \approx 16/9 \approx
1{\cdot}78$.

\paragraph{BDNK causal framework.}
The analysis above uses the BDNK (Bemfica--Disconzi--Noronha--Kovtun)
first-order causal dissipation framework
\citep{Bemfica2018,Kovtun2019}, in which the anisotropic stress is
given by the first-order constitutive relation
\begin{equation}
\label{eq:rel-7-BDNK}
\Pi^{\mu\nu} = -2\,\shearvisc\,\sigma^{\mu\nu}.
\end{equation}
Causality and strong hyperbolicity are ensured by constraints on the
BDNK frame coefficients rather than by relaxation equations.  In the
marginal stationary state ($\sigma_{\mathrm{rel}} = 0$) the
constitutive relation~\eqref{eq:rel-7-BDNK} is identical to what
one obtains from the Israel--Stewart theory at $\sigma=0$ (where
the relaxation terms vanish).  Therefore the critical Taylor number
for \emph{stationary} onset is the same in both formulations.

For the \emph{oscillatory} marginal state (examined in the
relativistic extension of~\S\,72), the BDNK and Israel--Stewart
theories differ: BDNK uses frequency-independent viscosity $\nu$
with frame-dependent causality bounds, while Israel--Stewart
introduces a frequency-dependent effective viscosity
$\nu/(1+\tau_{\pi}\sigma)$.  However, at the marginal state
$\sigma=0$ the two are identical.
